\section{$k$-positive maps, Schmidt rank, and Choi compression}

This section develops the $k$-positive-map criteria of
\cite[Chapter~3]{Wolf2012Quantum}. Schmidt-rank bounds convert positivity of
ampliations into quadratic-form conditions on the Choi matrix, while Ky
Fan's maximum principle controls the extremal overlap with vectors of bounded
Schmidt rank. The right-factor compression criteria give rectangular and
rank-$k$ projection forms of the same condition.

\subsection{Two-positive maps and the generalized Schwarz inequality}

\begin{definition}[$k$-positive map]\label{def:k_positive_map}
    \lean{IsNPositiveMap}
    \leanok
    A linear map $E : \MN{D} \to \MN{D}$ is $k$-positive if
    $E \otimes \id_{k}$ is positive on $\MN{D} \otimes \MN{k}$.
\end{definition}

\begin{definition}[Two-positive map]
    \label{def:2positive_map}
    \lean{Is2PositiveMap}
    \leanok
    A linear map is \emph{two-positive} if it is $2$-positive.
\end{definition}

\begin{theorem}[Completely positive implies two-positive]
    \label{thm:cp_implies_two_positive}
    \lean{IsCPMap.is2PositiveMap}
    \leanok
    \uses{def:cp_map, def:k_positive_map}
    Every completely positive map is two-positive.
\end{theorem}

\begin{proof}
    \leanok
    Let $E(X)=\sum_i K_iXK_i^\dagger$ be a Kraus representation. For every
    $Y\geq0$ in $\MN{D}\otimes\MN{2}$,
    \[
        (E\otimes\id_2)(Y)
        =\sum_i (K_i\otimes\Id_2)Y(K_i\otimes\Id_2)^\dagger\geq0,
    \]
    because each summand is positive semidefinite. Thus
    $E\otimes\id_2$ is positive, so $E$ is two-positive.
\end{proof}

\begin{theorem}[Two-positive implies positive]
    \label{thm:two_positive_implies_positive}
    \lean{Is2PositiveMap.isPositiveMap}
    \leanok
    \uses{def:positive_map, def:k_positive_map}
    Every two-positive map is positive.
\end{theorem}

\begin{proof}
    \leanok
    If $X\geq0$, then $X\otimes E_{11}\geq0$ in
    $\MN{D}\otimes\MN{2}$. Two-positivity gives
    \[
        (E\otimes\id_2)(X\otimes E_{11})=E(X)\otimes E_{11}\geq0.
    \]
    Its $(1,1)$ corner is $E(X)$, hence $E(X)\geq0$. Therefore $E$ is
    positive.
\end{proof}

\begin{definition}[Unital linear map]\label{def:unital_linear_map}
    \lean{KadisonSchwarz.IsUnitalMap}
    \leanok
    A linear map $E$ is unital if $E(\Id)=\Id$.
\end{definition}

\begin{theorem}[Kadison--Schwarz for unital two-positive maps]\label{thm:kadison_schwarz_2positive}
    \lean{kadison_schwarz_2positive}
    \leanok
    \uses{def:k_positive_map, def:unital_linear_map}
    If $E$ is unital and two-positive, then for all $X\in\MN{D}$,
    $E(X^\dagger X)\ge E(X)^\dagger E(X)$.
\end{theorem}

\begin{proof}
    \leanok
    Form the $2\times 2$ block matrix
    $(\begin{smallmatrix} \Id & X \\ X^\dagger & X^\dagger X \end{smallmatrix})\ge0$.
    Apply $E\otimes\id_{2}$ (positive by $2$-positivity) and take the Schur
    complement of the top-left identity block:
    $E(X^\dagger X)-E(X)^\dagger E(X)\ge0$.
\end{proof}

\begin{definition}[Blockwise ampliation]\label{def:blockwise_ampliation}
    \lean{nPositiveAmpliation}
    \leanok
    For a linear map $E : \MN{D}\to\MN{D}$, its $k$-fold ampliation acts on
    block matrices by
    \begin{align}
        E^{(k)}(X)_{(i,p),(j,q)}
         & =E((X_{(a,p),(b,q)})_{a,b})_{i,j}.
        \label{eq:schwarz_blockwise_ampliation}
    \end{align}
\end{definition}

\begin{theorem}[$k$-positivity as positivity of the ampliation]
    \lean{isNPositiveMap_iff_isPositiveMap_nPositiveAmpliation}
    \leanok
    \uses{def:k_positive_map, def:blockwise_ampliation}
    A linear map is $k$-positive if and only if its blockwise ampliation is
    positive on $\MN{D}\otimes\MN{k}$.
\end{theorem}

\begin{proof}
    \leanok
    This is precisely the preceding definition of the ampliation.
\end{proof}

\begin{definition}[Schmidt rank]\label{def:schmidt_rank}
    \lean{Matrix.schmidtCoeffMatrix, Matrix.schmidtRank, Matrix.HasSchmidtRankLE}
    \leanok
    A vector $\psi\in\C^{D}\otimes\C^{k}$ is identified with its coefficient
    matrix $C_\psi\in \MN{D,k}(\C)$. Its Schmidt rank is
    $\SR(\psi)=\rank(C_\psi)$. We write $\SR(\psi)\le r$ for the corresponding
    bounded-rank condition.
\end{definition}

\begin{theorem}[Basic bounds for Schmidt rank]\label{thm:schmidt_rank_basic_bounds}
    \lean{Matrix.schmidtRank_zero, Matrix.schmidtRank_le_left,
        Matrix.schmidtRank_le_right, Matrix.schmidtRank_le_min,
        Matrix.schmidtRank_product_le_one, Matrix.hasSchmidtRankLE_one_product}
    \leanok
    \uses{def:schmidt_rank}
    The zero vector has Schmidt rank zero, every vector satisfies
    $\SR(\psi)\le \min(D,k)$, and every product vector $u\otimes v$ has
    Schmidt rank at most one.
\end{theorem}

\begin{proof}
    \leanok
    The zero-vector claim is $\rank(0)=0$. For the dimension bounds,
    $\SR(\psi)=\rank(C_\psi)\le D$ and
    $\SR(\psi)=\rank(C_\psi)\le k$ follow from the row-rank and
    column-rank bounds on a $D\times k$ matrix, hence
    $\SR(\psi)\le \min(D,k)$. For a product vector, each column
    of $C_{u\otimes v}$ is a scalar multiple of $u$, so the column span has
    dimension at most one.
\end{proof}

\begin{definition}[Schmidt singular values]\label{def:schmidt_singular_values}
    \lean{Matrix.schmidtSingularValues}
    \leanok
    \uses{def:schmidt_rank}
    The Schmidt singular values of a vector
    $\psi\in\C^{D}\otimes\C^k$ are the singular values of the coefficient
    matrix $C_\psi$, viewed as a linear map $\C^k\to\C^D$.
\end{definition}

\begin{theorem}[Schmidt singular values and Schmidt rank]
    \label{thm:schmidt_singular_values_support}
    \lean{Matrix.support_schmidtSingularValues,
        Matrix.schmidtSingularValues_eq_zero_iff,
        Matrix.schmidtSingularValues_pos_iff_lt_schmidtRank,
        Matrix.hasSchmidtRankLE_iff_schmidtSingularValues_eq_zero}
    \leanok
    \uses{def:schmidt_rank, def:schmidt_singular_values}
    The Schmidt singular values are indexed from zero. The nonzero Schmidt
    singular values of $\psi$ are exactly those with index below
    $\SR(\psi)$. Equivalently, for zero-based indexing,
    $s_k(\psi)=0$ if and only if $\SR(\psi)\le k$.
\end{theorem}

\begin{proof}
    \leanok
    For a finite-dimensional linear map, the support of the singular-value
    sequence is the initial interval whose length is the dimension of the
    range. The coefficient matrix $C_\psi$ represents the corresponding map
    $\C^k\to\C^D$, and this range dimension is $\rank(C_\psi)$.
\end{proof}

\subsection{Schmidt rank and maximal overlap}

The monotonicity, rank-factorization, reduced-density, Rayleigh, and Parseval
steps used below are proved in
Section~\ref{sec:app_positive_not_cp_schmidt_rank}. The two extremal
projection bounds behind Ky Fan's maximum principle are proved in
Section~\ref{sec:app_positive_not_cp_ky_fan}.


\begin{definition}[Sum of the largest eigenvalues]
    \label{def:ky_fan_norm}
    \lean{Matrix.IsHermitian.kyFanNorm}
    \leanok
    Let $A\in\MN{D}$ be Hermitian with eigenvalues
    $\lambda_1\ge\lambda_2\ge\cdots\ge\lambda_D$ listed in decreasing order.
    For $0\le k\le D$, set
    \begin{align}
        S_k(A)
         & =\sum_{i=1}^{k}\lambda_i.
        \label{eq:channel_largest_eigenvalue_sum}
    \end{align}
    This is the sum of the $k$ largest signed eigenvalues. Indices beyond $D$
    contribute zero, so the value stabilizes at $\re\tr(A)$ once $k$ reaches
    $D$.

    The \emph{Ky-Fan $k$-norm} $\|A\|_{(k)}$ is by definition the sum of the
    $k$ largest \emph{singular values} of $A$. It agrees with $S_k(A)$ exactly
    when $A$ is positive semidefinite, where eigenvalues and singular values
    coincide; for an indefinite Hermitian $A$ the two differ. For
    $A=\diag(1,-2)$, one has $S_1(A)=1$ but $\|A\|_{(1)}=2$. The results below
    concern $S_k(A)$ for arbitrary Hermitian $A$, so on the
    positive-semidefinite cone they are statements about the Ky-Fan norm.
\end{definition}

\begin{theorem}[Ky Fan's maximum principle]
    \label{thm:ky_fan_maximum_principle}
    \lean{Matrix.IsHermitian.kyFanNorm_eq_sup_trace}
    \leanok
    \uses{def:ky_fan_norm, thm:ky_fan_achievability, thm:ky_fan_upper_bound}
    Let $A\in\MN{D}$ be Hermitian and $0\le k<D$. Then the sum of its $k$
    largest eigenvalues is the largest value of $\re\tr(PA)$ attained by an
    orthogonal projection $P$ of rank $k$:
    \begin{align}
        S_k(A)
         & =\max_{\substack{P^2=P=P^\dagger\\ \tr(P)=k}}\re\tr(PA).
        \label{eq:channel_ky_fan_maximum}
    \end{align}
    This is Ky Fan's maximum principle \cite{Fan1949Theorem}; see also
    \cite[Problem~I.6.15, Exercise~II.1.13]{Bhatia1997Matrix}. For a positive
    semidefinite $A$, where $S_k(A)$ coincides with the Ky-Fan norm
    $\|A\|_{(k)}$, it gives the extremal overlap of
    \cite[Lemma~3.1]{Wolf2012Quantum}.
\end{theorem}

\begin{proof}\leanok
    \uses{thm:ky_fan_achievability, thm:ky_fan_upper_bound}
    Theorem~\ref{thm:ky_fan_achievability} provides a rank-$k$ projection
    whose trace against $A$ equals $S_k(A)$, so the value is attained.
    Theorem~\ref{thm:ky_fan_upper_bound} shows that no rank-$k$ projection
    exceeds it. Hence $S_k(A)$ is the maximum.
\end{proof}

\begin{theorem}[Maximal overlap with a fixed Schmidt rank]
    \label{thm:maximal_overlap_schmidt_rank}
    \lean{Matrix.maximalSchmidtOverlap_eq_kyFanNorm}
    \leanok
    \uses{def:schmidt_rank, def:ky_fan_norm, thm:ky_fan_maximum_principle}
    Let $\phi\in\C^{D}\otimes\C^k$ be a normalized vector with reduced density
    matrix $\rho=\tr_2\ketbra{\phi}{\phi}$ on the first factor, and let
    $1\le n<D$. The largest squared overlap $|\braket{\phi}{\psi}|^2$ attained
    by a normalized vector $\psi$ of Schmidt rank at most $n$ is the Ky-Fan
    $n$-norm of $\rho$:
    \begin{align}
        \max_{\substack{\SR(\psi)\le n\\ \|\psi\|=1}}
        |\braket{\phi}{\psi}|^2
         & =\|\rho\|_{(n)}.
        \label{eq:schwarz_maximal_schmidt_overlap}
    \end{align}
    Because $\rho$ is positive semidefinite, this norm is the sum of its $n$
    largest eigenvalues, which coincide there with its singular values. This is
    \cite[Lemma~3.1]{Wolf2012Quantum} for $1\le n<D$; the top index $n=D$,
    where the value is $\tr\rho=1$, is
    Theorem~\ref{thm:maximal_overlap_schmidt_rank_top}.
\end{theorem}

\begin{proof}
    \leanok
    \uses{thm:ky_fan_maximum_principle}
    Write $C$ for the coefficient matrix of $\phi$, so that
    $\rho=CC^\dagger$ and the overlap with a vector $\psi$ of coefficient
    matrix $B$ is the Frobenius pairing
    $\braket{\phi}{\psi}=\tr(C^\dagger B)$; normalization reads
    $\tr(B^\dagger B)=1$ and the Schmidt-rank constraint reads
    $\rank(B)\le n$.

    For the upper bound, let $P$ be the orthogonal projection onto the column
    space of $B$, so $\rank(P)\le n$ and $PB=B$. The Cauchy--Schwarz inequality
    for the Frobenius pairing gives
    \begin{align}
        |\tr(C^\dagger B)|^2
         & =|\tr((PC)^\dagger B)|^2
        \le \tr\!((PC)^\dagger(PC))\tr(B^\dagger B)
        =\tr(P\rho).
        \notag
    \end{align}
    Since $P$ has rank at most $n$ and $\rho$ is positive semidefinite, the
    positive-semidefinite form of Ky Fan's maximum principle
    (Theorem~\ref{thm:ky_fan_maximum_principle}) gives
    $\tr(P\rho)\le \|\rho\|_{(n)}$. Thus
    $|\braket{\phi}{\psi}|^2\le\|\rho\|_{(n)}$ for every admissible $\psi$.

    For the converse, let $P_n$ be the eigenprojection of $\rho$ onto its $n$
    largest eigenvalues, so $\tr(P_n\rho)=\|\rho\|_{(n)}$, and set
    \begin{align}
        B
         & =\|\rho\|_{(n)}^{-1/2}P_nC.
        \label{eq:schwarz_schmidt_overlap_witness}
    \end{align}
    Then $\rank(B)\le\rank(P_n)=n$ and
    $\tr(B^\dagger B)=\|\rho\|_{(n)}^{-1}\tr(P_n\rho)=1$, while
    \begin{align}
        |\tr(C^\dagger B)|^2
         & =\|\rho\|_{(n)}^{-1}|\tr(P_n\rho)|^2
        =\|\rho\|_{(n)}.
        \notag
    \end{align}
    Thus the bound is attained. For a normalized $\phi$ with $1\le n$, one
    has $\|\rho\|_{(n)}\ge\|\rho\|_{(1)}>0$, since $\tr\rho=1$, which
    justifies the normalization in
    (\ref{eq:schwarz_schmidt_overlap_witness}).
\end{proof}

\begin{theorem}[Maximal overlap at the top Schmidt rank]
    \label{thm:maximal_overlap_schmidt_rank_top}
    \lean{Matrix.maximalSchmidtOverlap_eq_kyFanNorm_top}
    \leanok
    \uses{def:schmidt_rank, def:ky_fan_norm}
    At the top index $n=D$, the Schmidt-rank constraint $\SR(\psi)\le D$
    is vacuous, so the largest squared overlap is the Ky-Fan $D$-norm of
    $\rho$:
    \begin{align}
        \max_{\|\psi\|=1}|\braket{\phi}{\psi}|^2
         & =\|\rho\|_{(D)}.
        \label{eq:schwarz_maximal_overlap_top}
    \end{align}
    The maximum is attained at $\psi=\phi$. The Ky-Fan $D$-norm sums all
    eigenvalues, so for a normalized $\phi$ it reduces to
    $\|\rho\|_{(D)}=\tr\rho=1$. Together with
    Theorem~\ref{thm:maximal_overlap_schmidt_rank}, this covers the full range
    $1\le n\le D$ of \cite[Lemma~3.1]{Wolf2012Quantum}.
\end{theorem}

\begin{proof}
    \leanok
    The Ky-Fan $D$-norm sums all eigenvalues of $\rho$, so
    $\|\rho\|_{(D)}=\tr\rho=\|\phi\|^2=1$. The bound $1$ is attained at
    $\psi=\phi$, whose Schmidt rank is at most $D$; for any normalized
    $\psi$, the Cauchy--Schwarz inequality gives
    $|\braket{\phi}{\psi}|^2\le\|\phi\|^2\|\psi\|^2=1$.
\end{proof}

\begin{definition}[Reduced density of an eigenvector]
    \label{def:reduced_eig_density}
    \lean{Matrix.IsHermitian.reducedEigDensity}
    \leanok
    Let $\tau$ be a Hermitian operator on
    $\C^{D'}\otimes\C^{D}$ with normalized eigenvectors $\phi_i$. The reduced
    density operator of the $i$-th eigenvector on the first factor is
    $\rho_i=\tr_2\ketbra{\phi_i}{\phi_i}$.
\end{definition}

\begin{theorem}[Spectral lower bound for the Schmidt-rank expectation]
    \label{thm:spectral_lower_bound_schmidt_rank}
    \lean{Matrix.IsHermitian.spectral_lower_bound}
    \leanok
    \uses{def:schmidt_rank, def:ky_fan_norm, def:reduced_eig_density,
        lem:reduced_eig_density_posSemidef, thm:maximal_overlap_schmidt_rank}
    Let $\tau$ be a Hermitian operator on $\C^{D'}\otimes\C^{D}$ with
    eigenvalues $\nu_i$ and normalized eigenvectors $\phi_i$, and write
    $\rho_i=\tr_2\ketbra{\phi_i}{\phi_i}$ for the reduced density operator of
    $\phi_i$ on the first factor. Let $1\le n<D'$ and let $\nu_0\ge 0$ be a
    lower bound for the positive eigenvalues (the smallest positive
    eigenvalue). Then every normalized vector $\psi$ of Schmidt rank at most
    $n$ satisfies
    \begin{align}
        \nu_0+\sum_{i:\nu_i\le0}
        (\nu_i-\nu_0)\|\rho_i\|_{(n)}
         & \le\bra{\psi}\tau\ket{\psi}.
        \label{eq:schwarz_spectral_schmidt_lower}
    \end{align}
    Hence the same bound holds for the infimum over such vectors. This is the
    lower bound of \cite[Chapter~3, Proposition~3.2]{Wolf2012Quantum} for
    $n<D'$; the top index $n=D'$ is
    Theorem~\ref{thm:spectral_bound_schmidt_rank_top}.
\end{theorem}

\begin{proof}
    \leanok
    \uses{lem:eigenbasis_rayleigh, lem:eigenbasis_parseval,
        thm:maximal_overlap_schmidt_rank}
    By (\ref{eq:schwarz_eigenbasis_rayleigh}),
    $\bra{\psi}\tau\ket{\psi}
        =\sum_i\nu_i|\braket{\phi_i}{\psi}|^2$, while
    (\ref{eq:schwarz_eigenbasis_parseval}) gives
    $\sum_i|\braket{\phi_i}{\psi}|^2=\|\psi\|^2=1$. Subtracting $\nu_0$
    times the latter from the former gives
    \begin{align}
        \bra{\psi}\tau\ket{\psi}
         & =\nu_0+\sum_i(\nu_i-\nu_0)
        |\braket{\phi_i}{\psi}|^2.
        \label{eq:schwarz_spectral_lower_expansion}
    \end{align}
    For an eigenvalue $\nu_i>0$, the factor $\nu_i-\nu_0$ is non-negative, so
    those terms only increase the sum and may be dropped. For an eigenvalue
    $\nu_i\le0$, the factor $\nu_i-\nu_0$ is nonpositive, and
    Theorem~\ref{thm:maximal_overlap_schmidt_rank} gives
    $|\braket{\phi_i}{\psi}|^2\le\|\rho_i\|_{(n)}$. Multiplication by the
    nonpositive factor reverses the inequality, which gives
    (\ref{eq:schwarz_spectral_schmidt_lower}).
\end{proof}

\begin{theorem}[Spectral upper bound for the Schmidt-rank expectation]
    \label{thm:spectral_upper_bound_schmidt_rank}
    \lean{Matrix.IsHermitian.exists_le_spectral_upper_bound}
    \leanok
    \uses{def:schmidt_rank, def:ky_fan_norm, def:reduced_eig_density,
        lem:reduced_eig_density_posSemidef, thm:maximal_overlap_schmidt_rank}
    With $\tau$, $\phi_i$, and $\rho_i$ as above, let $1\le n<D'$ and suppose
    every eigenvalue of $\tau$ is at most $\nu$. For any eigenvector index $j$,
    there is a normalized vector $\psi$ of Schmidt rank at most $n$ with
    \begin{align}
        \bra{\psi}\tau\ket{\psi}
         & \le\nu+(\nu_j-\nu)\|\rho_j\|_{(n)}.
        \label{eq:schwarz_spectral_schmidt_upper}
    \end{align}
    Hence the infimum over such vectors lies below that value. Taking $\nu$
    to be the largest positive eigenvalue and $j$ the index of the unique
    non-positive eigenvalue $\nu_-$ recovers the upper bound
    $\nu+(\nu_--\nu)\|\rho_-\|_{(n)}$ of
    \cite[Chapter~3, Proposition~3.2]{Wolf2012Quantum}. The statement is for
    $n<D'$; the top index $n=D'$ is
    Theorem~\ref{thm:spectral_bound_schmidt_rank_top}.
\end{theorem}

\begin{proof}
    \leanok
    \uses{lem:eigenbasis_rayleigh, lem:eigenbasis_parseval,
        thm:maximal_overlap_schmidt_rank}
    Replacing $\nu_0$ by $\nu$ in
    (\ref{eq:schwarz_spectral_lower_expansion}) gives
    \begin{align}
        \bra{\psi}\tau\ket{\psi}
         & =\nu+\sum_i(\nu_i-\nu)|\braket{\phi_i}{\psi}|^2.
        \notag
    \end{align}
    Every factor $\nu_i-\nu$ is nonpositive, so all terms are nonpositive and
    keeping only the $i=j$ term gives
    $\bra{\psi}\tau\ket{\psi}
        \le\nu+(\nu_j-\nu)|\braket{\phi_j}{\psi}|^2$.
    Theorem~\ref{thm:maximal_overlap_schmidt_rank} supplies a normalized
    $\psi$ of Schmidt rank at most $n$ attaining
    $|\braket{\phi_j}{\psi}|^2=\|\rho_j\|_{(n)}$; at that vector the
    right-hand side is the expression in
    (\ref{eq:schwarz_spectral_schmidt_upper}).
\end{proof}

\begin{theorem}[Spectral criterion at the top Schmidt rank]
    \label{thm:spectral_bound_schmidt_rank_top}
    \lean{Matrix.IsHermitian.spectral_lower_bound_top,
        Matrix.IsHermitian.exists_spectral_lower_bound_top}
    \leanok
    \uses{lem:eigenbasis_rayleigh, lem:eigenbasis_parseval}
    At the top index $n=D'$, the Schmidt-rank constraint is vacuous and each
    Ky-Fan $D'$-norm reduces to $\|\rho_i\|_{(D')}=\tr\rho_i=1$. Thus the
    lower and upper bounds of
    \cite[Chapter~3, Proposition~3.2]{Wolf2012Quantum} reduce to the Rayleigh
    characterization of the least eigenvalue: the lower bound becomes
    $\nu_{\min}\le\bra{\psi}\tau\ket{\psi}$ for every normalized $\psi$, and
    the upper bound on the infimum becomes the existence of a normalized
    eigenvector with
    $\bra{\psi}\tau\ket{\psi}=\nu_{\min}$. Together they give
    \begin{align}
        \inf_{\|\psi\|=1}\bra{\psi}\tau\ket{\psi}
         & =\nu_{\min}.
        \label{eq:schwarz_spectral_schmidt_top}
    \end{align}
    There is no Schmidt-rank restriction. Together with
    Theorems~\ref{thm:spectral_lower_bound_schmidt_rank}
    and~\ref{thm:spectral_upper_bound_schmidt_rank}, this covers the full
    range $1\le n\le D'$ of
    \cite[Chapter~3, Proposition~3.2]{Wolf2012Quantum}.
\end{theorem}

\begin{proof}
    \leanok
    \uses{lem:eigenbasis_rayleigh, lem:eigenbasis_parseval}
    The eigenbasis expansion (\ref{eq:schwarz_eigenbasis_rayleigh}), together
    with Parseval's identity (\ref{eq:schwarz_eigenbasis_parseval}), bounds
    each weight $\nu_i$ below by $\nu_{\min}$, giving
    $\nu_{\min}\le\bra{\psi}\tau\ket{\psi}$. Evaluating at the eigenvector
    $\phi_j$ for a least-eigenvalue index $j$ gives
    $\bra{\phi_j}\tau\ket{\phi_j}=\nu_j=\nu_{\min}$, so the lower bound is
    attained and the infimum is (\ref{eq:schwarz_spectral_schmidt_top}).
\end{proof}

\begin{theorem}[Rank-one test for $k$-positivity]\label{thm:k_positive_rank_one_test}
    \lean{isNPositiveMap_iff_forall_ampliation_rank_one_posSemidef}
    \leanok
    \uses{def:k_positive_map, def:blockwise_ampliation}
    A linear map $E:\MN{D}\to\MN{D}$ is $k$-positive if and only if, for
    every vector $\phi\in\C^D\otimes\C^k$,
    $E^{(k)}(\ketbra{\phi}{\phi})\ge0$.
    This is the pure-state reduction used in
    \cite[Chapter~3, Proposition~3.1]{Wolf2012Quantum}.
\end{theorem}

\begin{proof}
    \leanok
    The forward implication is the definition of $k$-positivity, since
    $\ketbra{\phi}{\phi}$ is positive semidefinite. Conversely, every
    positive semidefinite matrix on $\C^D\otimes\C^k$ is a finite sum
    $\sum_i\ketbra{\phi_i}{\phi_i}$, and
    \begin{align}
        E^{(k)}\!\left(\sum_i\ketbra{\phi_i}{\phi_i}\right)
         & =\sum_iE^{(k)}(\ketbra{\phi_i}{\phi_i})\ge0.
        \notag
    \end{align}
    by linearity of $E^{(k)}$, the hypothesis, and closure of the positive
    semidefinite cone under finite sums.
\end{proof}

\subsection{Choi compression criteria}

The right-tensor identities, parametrizations of bounded Schmidt rank, and
projection factorizations used in this subsection are proved in
Section~\ref{sec:app_positive_not_cp_choi_compression}.

\begin{definition}[Right compression of the Choi matrix]\label{def:choi_right_compression}
    \lean{ChoiJamiolkowski.rightCompression,
        ChoiJamiolkowski.rightCompression_apply,
        ChoiJamiolkowski.compressedOmegaVector}
    \leanok
    \uses{def:blockwise_ampliation}
    For $X\in\MN{D,k}(\C)$ and a Choi matrix $\tau$, the right-factor
    compression in the blockwise ampliation index convention has entries
    \begin{align}
        C_{(i,p),(j,q)}
         & =\sum_{a,b}X_{a,p}\tau_{(i,a),(j,b)}
        \overline{X_{b,q}}.
        \label{eq:schwarz_right_choi_compression}
    \end{align}
    The associated vector has coefficients $D^{-1/2}X_{i,p}$.
\end{definition}

\begin{definition}[Right tensor factor for Choi compression]
    \label{def:choi_right_tensor_factor}
    \lean{ChoiJamiolkowski.rightTensorMatrix}
    \leanok
    For $X\in\MN{D,k}(\C)$, let $R_X$ be the matrix from
    $\C^D\otimes\C^D$ to $\C^D\otimes\C^k$ with entries
    \begin{align}
        (R_X)_{(i,p),(j,a)}
         & =\delta_{i,j}X_{a,p}.
        \label{eq:schwarz_right_tensor_factor}
    \end{align}
\end{definition}
\begin{theorem}[Schmidt-rank Choi expectation from $k$-positivity]
    \label{thm:k_positive_schmidt_rank_choi_expectation}
    \lean{ChoiJamiolkowski.IsNPositiveMap.choiMatrix_quadraticForm_nonneg_of_hasSchmidtRankLE}
    \leanok
    \uses{lem:right_tensor_adjoint_schmidt_rank_representatives,
        thm:choi_right_compression_quadratic_form,
        thm:k_positive_iff_right_choi_compressions}
    Assume $D>0$. If $T$ is $k$-positive, then
    $\langle\psi,\tau_T\psi\rangle\ge0$ for every
    $\psi\in\C^D\otimes\C^D$ with $\SR(\psi)\le k$.
\end{theorem}

\begin{proof}
    \leanok
    Write $\psi=R_X^\dagger\eta$ using
    Lemma~\ref{lem:right_tensor_adjoint_schmidt_rank_representatives}. By
    Theorem~\ref{thm:k_positive_iff_right_choi_compressions}, the
    compression $C_T(X)$ is positive semidefinite. Its quadratic form at
    $\eta$ is therefore non-negative, and
    (\ref{eq:schwarz_right_compression_quadratic}) identifies this number
    with $\langle\psi,\tau_T\psi\rangle$.
\end{proof}

\begin{theorem}[Converse Schmidt-rank Choi test]
    \label{thm:schmidt_rank_choi_test_converse}
    \lean{ChoiJamiolkowski.isNPositiveMap_of_forall_hasSchmidtRankLE_choiMatrix_quadraticForm_nonneg}
    \leanok
    \uses{lem:right_tensor_adjoint_schmidt_rank,
        thm:choi_right_compression_quadratic_form,
        thm:k_positive_iff_right_choi_compressions}
    Assume $D>0$. If $\langle\psi,\tau_T\psi\rangle\ge0$ for every
    $\psi\in\C^D\otimes\C^D$ with $\SR(\psi)\le k$, then $T$ is
    $k$-positive.
\end{theorem}

\begin{proof}
    \leanok
    By Theorem~\ref{thm:k_positive_iff_right_choi_compressions}, it is enough
    to prove that $C_T(X)$ is positive semidefinite for each
    $X\in\MN{D,k}(\C)$. For any $\eta$,
    (\ref{eq:schwarz_right_compression_quadratic}) identifies
    $\langle\eta,C_T(X)\eta\rangle$ with
    $\langle R_X^\dagger\eta,\tau_T R_X^\dagger\eta\rangle$, and
    Lemma~\ref{lem:right_tensor_adjoint_schmidt_rank} gives
    $\SR(R_X^\dagger\eta)\le k$. Thus every quadratic form of $C_T(X)$ is
    non-negative. Over a complex finite-dimensional space, polarization
    recovers Hermiticity from this real non-negative quadratic-form condition,
    so $C_T(X)$ is positive semidefinite.
\end{proof}

\begin{theorem}[Schmidt-rank Choi criterion]
    \label{thm:schmidt_rank_choi_test_iff}
    \lean{ChoiJamiolkowski.isNPositiveMap_iff_forall_hasSchmidtRankLE_choiMatrix_quadraticForm_nonneg}
    \leanok
    \uses{thm:k_positive_schmidt_rank_choi_expectation,
        thm:schmidt_rank_choi_test_converse}
    Assume $D>0$. Then $T$ is $k$-positive if and only if
    $\langle\psi,\tau_T\psi\rangle\ge0$ for every
    $\psi\in\C^D\otimes\C^D$ with $\SR(\psi)\le k$.
\end{theorem}

\begin{proof}
    \leanok
    Combine Theorem~\ref{thm:k_positive_schmidt_rank_choi_expectation} with
    Theorem~\ref{thm:schmidt_rank_choi_test_converse}.
\end{proof}
\begin{theorem}[Rectangular Choi-compression test for $k$-positivity]
    \label{thm:k_positive_iff_right_choi_compressions}
    \lean{ChoiJamiolkowski.isNPositiveMap_iff_forall_rightCompression_posSemidef}
    \leanok
    \uses{thm:k_positive_rank_one_test, def:choi_right_compression,
        lem:rank_one_ampliation_choi_compression}
    Assume $D>0$. A linear map
    $T:\MN{D,D}(\C)\to\MN{D,D}(\C)$ is $k$-positive if and only if, for every
    $X\in\MN{D,k}(\C)$, the right-factor compression of the Choi matrix of
    $T$ by $X$ is positive semidefinite.
\end{theorem}

\begin{proof}
    \leanok
    The forward direction applies the rank-one $k$-positivity test to the
    vector $\psi_X$ in (\ref{eq:schwarz_compressed_omega}), then uses
    Lemma~\ref{lem:rank_one_ampliation_choi_compression} in the form
    \begin{align}
        (T\otimes\id_k)(\ketbra{\psi_X}{\psi_X})
         & =C_T(X).
        \label{eq:schwarz_ampliation_compression}
    \end{align}
    Conversely, given a vector $\psi$ on the ampliated space, take
    $X_{i,p}=D^{1/2}\psi_{(i,p)}$. Since $D>0$, this choice satisfies
    $D^{-1/2}X_{i,p}=\psi_{(i,p)}$. Substituting this vector in
    (\ref{eq:schwarz_ampliation_compression}) gives
    $(T\otimes\id_k)(\ketbra{\psi}{\psi})=C_T(X)$.
    Positivity of $C_T(X)$ is therefore positivity of the ampliation on the
    rank-one matrix $\ketbra{\psi}{\psi}$, and the rank-one test gives
    $k$-positivity.
\end{proof}

\begin{theorem}[Right tensor Choi sandwiches of a $k$-positive map]
    \label{thm:k_positive_right_tensor_choi_sandwiches}
    \lean{ChoiJamiolkowski.IsNPositiveMap.rightTensor_choiMatrix_sandwich_posSemidef}
    \leanok
    \uses{thm:k_positive_iff_right_choi_compressions,
        thm:choi_right_tensor_sandwich}
    Assume $D>0$. If
    $T:\MN{D,D}(\C)\to\MN{D,D}(\C)$ is $k$-positive, then, for every
    $X\in\MN{D,k}(\C)$,
    \begin{align}
        R_X\tau R_X^\dagger
         & \ge0,
        \label{eq:schwarz_k_positive_choi_sandwich}
    \end{align}
    where $\tau$ is the Choi matrix of $T$.
\end{theorem}

\begin{proof}
    \leanok
    By Theorem~\ref{thm:k_positive_iff_right_choi_compressions},
    $C_T(X)\ge0$ for every $X\in\MN{D,k}(\C)$. The sandwich formula
    (\ref{eq:schwarz_right_tensor_sandwich}) identifies
    $R_X\tau R_X^\dagger$ with $C_T(X)$.
\end{proof}
\begin{theorem}[Rank-$k$ right projections in the Choi criterion]
    \label{thm:rank_k_right_projection_choi_sandwich_forward}
    \lean{ChoiJamiolkowski.IsNPositiveMap.rightTensor_choiMatrix_sandwich_posSemidef_of_isHermitian_idempotent_rank}
    \leanok
    \uses{lem:rank_k_hermitian_projection_factorization,
        thm:factorized_right_projection_choi_sandwich}
    Assume $D>0$. Let $T:\MN{D,D}(\C)\to\MN{D,D}(\C)$ be $k$-positive,
    with Choi matrix $\tau$. If $P\in\MN{D,D}(\C)$ is Hermitian, satisfies
    $P^2=P$, and has rank $k$, then $R_P\tau R_P^\dagger\geq0$.
    This is the forward rank-$k$ projection-compression direction of
    \cite[Chapter~3, Proposition~3.1]{Wolf2012Quantum}.
\end{theorem}

\begin{proof}
    \leanok
    By Lemma~\ref{lem:rank_k_hermitian_projection_factorization}, write
    $P=VV^\dagger$ with $V\in\MN{D,k}(\C)$. The conclusion is then exactly
    Theorem~\ref{thm:factorized_right_projection_choi_sandwich}.
\end{proof}

\begin{theorem}[Projection-compression converse for the Choi criterion]
    \label{thm:rank_k_right_projection_choi_sandwich_converse}
    \lean{ChoiJamiolkowski.isNPositiveMap_of_forall_rankProjection_rightTensor_choiMatrix_sandwich_posSemidef}
    \leanok
    \uses{thm:k_positive_iff_right_choi_compressions,
        lem:square_sandwich_fixed_rectangular_compression,
        lem:rank_k_projection_fixing_rectangular_factor}
    Assume $D>0$ and $k\le D$. Let
    $T:\MN{D,D}(\C)\to\MN{D,D}(\C)$ have Choi matrix $\tau$. Suppose that
    $R_P\tau R_P^\dagger\geq0$ for every Hermitian projection
    $P\in\MN{D,D}(\C)$ of rank $k$. Then $T$ is $k$-positive.
\end{theorem}

\begin{proof}
    \leanok
    By Theorem~\ref{thm:k_positive_iff_right_choi_compressions}, it is enough
    to prove positivity of the right-factor compression by an arbitrary
    $X\in\MN{D,k}(\C)$. Choose a rank-$k$ Hermitian projection $P$ with
    $PX=X$ using
    Lemma~\ref{lem:rank_k_projection_fixing_rectangular_factor}. The assumed
    positivity of $R_P\tau R_P^\dagger$, together with
    Lemma~\ref{lem:square_sandwich_fixed_rectangular_compression}, gives
    positivity of the rectangular compression by $X$.
\end{proof}

\begin{theorem}[Rank-$k$ projection form of the Choi criterion]
    \label{thm:rank_k_right_projection_choi_criterion}
    \lean{ChoiJamiolkowski.isNPositiveMap_iff_forall_rankProjection_rightTensor_choiMatrix_sandwich_posSemidef}
    \leanok
    \uses{thm:rank_k_right_projection_choi_sandwich_forward,
        thm:rank_k_right_projection_choi_sandwich_converse}
    Assume $D>0$ and $k\le D$. A linear map
    $T:\MN{D,D}(\C)\to\MN{D,D}(\C)$ is $k$-positive if and only if, for every
    Hermitian projection $P\in\MN{D,D}(\C)$ of rank $k$,
    $R_P\tau R_P^\dagger\geq0$, where $\tau$ is the Choi matrix of $T$.
    This is the rank-$k$ projection formulation of
    \cite[Chapter~3, Proposition~3.1]{Wolf2012Quantum}.
\end{theorem}

\begin{proof}
    \leanok
    The forward direction is
    Theorem~\ref{thm:rank_k_right_projection_choi_sandwich_forward}. The
    reverse direction is
    Theorem~\ref{thm:rank_k_right_projection_choi_sandwich_converse}.
\end{proof}

\section{Trace adjoints and the positivity hierarchy}

The trace-pairing adjoint was introduced in
Definition~\ref{def:trace_pairing_adjoint}. Its compatibility with
ampliation makes $k$-positivity invariant under trace adjoints and identifies
the endpoints of the positivity hierarchy. Routine closure properties of the
cones of $k$-positive maps are collected in
Section~\ref{sec:app_positive_not_cp_k_positive_closure}.

\begin{theorem}[Trace-pairing adjoint of an ampliation]
    \label{thm:trace_pairing_adjoint_ampliation}
    \lean{nPositiveAmpliation_traceAdjointMap}
    \leanok
    \uses{def:blockwise_ampliation, def:trace_pairing_adjoint,
        thm:trace_pairing_adjoint_identity}
    For every linear map $E : \MN{D}\to\MN{D}$, the trace-pairing adjoint
    commutes with ampliation:
    \begin{align}
        (E^{(k)})^*
         & =(E^*)^{(k)}.
        \label{eq:schwarz_trace_adjoint_ampliation}
    \end{align}
\end{theorem}

\begin{proof}
    \leanok
    For block matrices write $\rho_{pq}$ and $X_{pq}$ for their
    $\MN{D}$-valued blocks. Expanding the trace by blocks and applying
    (\ref{eq:schwarz_trace_adjoint_pairing}) gives
    \begin{align}
        \tr((E^*)^{(k)}(\rho)X)
         & =\sum_{p,q}\tr(E^*(\rho_{pq})X_{qp})
        =\sum_{p,q}\tr(\rho_{pq}E(X_{qp}))
        =\tr(\rho E^{(k)}(X)).
        \notag
    \end{align}
    Nondegeneracy of the trace pairing identifies the two adjoints, proving
    (\ref{eq:schwarz_trace_adjoint_ampliation}).
\end{proof}

\begin{theorem}[Trace-pairing adjoint of a $k$-positive map]
    \label{thm:k_positive_trace_pairing_adjoint}
    \lean{IsNPositiveMap.traceAdjointMap}
    \leanok
    \uses{def:k_positive_map, thm:trace_pairing_adjoint_ampliation,
        thm:trace_pairing_adjoint_positive}
    If $E$ is $k$-positive, then $E^*$ is $k$-positive.
\end{theorem}

\begin{proof}
    \leanok
    Since $E$ is $k$-positive, its ampliation $E^{(k)}$ is positive. The
    trace-pairing adjoint of this positive map is positive, and
    (\ref{eq:schwarz_trace_adjoint_ampliation}) identifies that adjoint with
    $(E^*)^{(k)}$.
\end{proof}

\begin{theorem}[$k$-positivity is trace-adjoint invariant]
    \label{thm:k_positive_trace_pairing_adjoint_iff}
    \lean{isNPositiveMap_traceAdjointMap_iff}
    \leanok
    \uses{thm:k_positive_trace_pairing_adjoint,
        thm:trace_pairing_adjoint_involutive}
    A linear map is $k$-positive if and only if its trace-pairing adjoint is
    $k$-positive.
\end{theorem}

\begin{proof}
    \leanok
    One direction is
    Theorem~\ref{thm:k_positive_trace_pairing_adjoint}. The other follows by
    applying the same theorem to $E^*$ and using
    Theorem~\ref{thm:trace_pairing_adjoint_involutive}.
\end{proof}

\begin{theorem}[Completely positive maps are $k$-positive]
    \label{thm:cp_implies_k_positive}
    \lean{IsCPMap.isNPositiveMap}
    \leanok
    \uses{def:cp_map}
    Every completely positive map is $k$-positive for all $k \ge 1$.
\end{theorem}

\begin{proof}
    \leanok
    Choose a Kraus representation $E(X)=\sum_i K_iXK_i^\dagger$. For every
    $Y\succeq0$ in $\MN{D}\otimes\MN{k}$,
    \begin{align}
        (E\otimes\id_k)(Y)
         & =\sum_i (K_i\otimes\Id_k)Y(K_i\otimes\Id_k)^\dagger
        \succeq0.
        \notag
    \end{align}
    Thus every ampliation $E\otimes\id_k$ is positive.
\end{proof}

\begin{theorem}[One-positive maps are positive]
    \label{thm:one_positive_iff_positive}
    \lean{isNPositiveMap_one_iff_isPositiveMap}
    \leanok
    \uses{def:k_positive_map, def:positive_map}
    A linear map is $1$-positive if and only if it is positive.
\end{theorem}

\begin{proof}
    \leanok
    The ampliation by $\MN{1}$ is identified with the original matrix algebra:
    all block indices have the unique value, so positivity of
    $E\otimes\id_1$ is exactly positivity of $E$.
\end{proof}

\begin{theorem}[$D$-positive maps on nonzero $\MN{D}$ are completely positive]
    \label{thm:d_positive_iff_completely_positive}
    \lean{isCPMap_iff_isNPositiveMap_card}
    \leanok
    \uses{def:k_positive_map, def:cp_map, thm:cp_iff_choi_posSemidef}
    For maps $E:\MN{D}\to\MN{D}$ with $D\ge 1$, complete positivity is
    equivalent to $D$-positivity.
\end{theorem}

\begin{proof}
    \leanok
    Complete positivity implies $k$-positivity for every $k$. Conversely,
    assume $D\ge 1$. If $E$ is $D$-positive, apply $E\otimes\id_D$ to the
    maximally entangled projector $\ketbra{\Omega}{\Omega}$. The result is
    the Choi matrix of $E$, hence it is positive semidefinite. Choi's theorem
    then gives complete positivity.
\end{proof}
